\vspace{-2.5em}
\section{Model Framework}
\vspace{-1em}
\subsection{Setup}
\vspace{-0.5em}

The household consists of two members $j \in \{1,2\}$ over a finite working life $t = 1, \dots, T$.
Each period they jointly choose consumption, hours of work, and savings.
The state variables are household wealth $A_t$ and each member's human capital stock $k_{j,t}$.
\vspace{-1em}
\subsection{Household Problem}

\vspace{-0.5em}
\paragraph{Wages.} Earnings capacity depends on accumulated human capital:
\vspace{-0.5em}
\begin{equation}
    w_{j,t} = \exp\!\bigl(\bar{w}_j + \alpha_j \, k_{j,t}\bigr),
\end{equation}
\vspace{-0.5em}
where $\bar{w}_j$ is a gender-specific wage intercept and $\alpha_j \geq 0$ is the return to human capital.

\paragraph{Human capital.} Hours worked serve as on-the-job investment:
\vspace{-0.5em}
\begin{equation}
    k_{j,t+1} = (1-\delta)\,k_{j,t} + h_{j,t}, \qquad \delta \in (0,1).
\end{equation}
\vspace{-0.5em}
\noindent \textbf{Preferences.} Each member has instantaneous utility over private consumption $c_{j,t}$ and hours $h_{j,t} \in [0,1]$:
\begin{equation}
    u_j(c_{j,t},\, h_{j,t}) \;=\; \frac{c_{j,t}^{1+\eta}}{1+\eta} \;-\; \rho_j\,\frac{h_{j,t}^{1+\gamma}}{1+\gamma},
\end{equation}
where $\rho_j$ captures gender-specific disutility of labor and $\gamma > 0$ governs the Frisch elasticity. The parameter $\eta < -1$ follows the course convention (e.g.\ $\eta = -1.5$ in Lecture 06), so that $1 + \eta < 0$ and utility is increasing and concave in consumption.
\vspace{-0.5em}
\paragraph{Budget constraint.} Total consumption plus savings equals after-tax household income:
\vspace{-0.5em}
\begin{equation}
    c_{1,t} + c_{2,t} + A_{t+1} \;=\; (1+r)\,A_t \;+\; Y^{\text{net}}_t,
\end{equation}
\vspace{-0.5em}
where after-tax income depends on the tax regime:
\vspace{-0.5em}
\begin{align}
    Y^{\text{joint}}_t  &= \bigl[1 - \tau(w_{1,t}h_{1,t} + w_{2,t}h_{2,t})\bigr]\,(w_{1,t}h_{1,t} + w_{2,t}h_{2,t}), \\
    Y^{\text{indiv}}_t  &= \bigl[1-\tau(w_{1,t}h_{1,t})\bigr]\,w_{1,t}h_{1,t} \;+\; \bigl[1-\tau(w_{2,t}h_{2,t})\bigr]\,w_{2,t}h_{2,t}.
\end{align}
Here $\tau(\cdot)$ is a progressive marginal tax schedule.
Under joint taxation the secondary earner's income enters the household total and is taxed at the household's top bracket; under individual taxation each earner faces their own marginal rate.
\vspace{-0.5em}
\paragraph{Value function.} The household maximizes discounted joint utility with equal Pareto weights:
\begin{equation}
    V_t(A_t, k_{1,t}, k_{2,t}) \;=\; \max_{c_1,\, c_2,\, h_1,\, h_2,\, A_{t+1}}
    \Bigl[\, u_1 + u_2 \;+\; \beta\, V_{t+1}(A_{t+1}, k_{1,t+1}, k_{2,t+1}) \,\Bigr],
\end{equation}
solved by backward induction from the terminal period $T$. Equal weights are assumed for tractability; an natural extension would endogenize the Pareto weight through Nash bargaining, as in Lectures 08--09, which could amplify or dampen the tax reform effect depending on how the reform shifts the spouses' outside options.
\vspace{-0.5em}
\subsection{Equilibrium}

The tax schedule $\tau(\cdot)$ is taken as exogenous and regime-specific.
A revenue-neutral comparison sets the individual-tax schedule to raise the same total revenue as the joint-tax baseline, ensuring that any behavioral differences reflect the structure of taxation rather than its level.
The key mechanism is that switching to individual taxation lowers the effective marginal rate on the secondary earner, raising that member's labor supply and --- through equation~(2) --- their human capital and future wages, amplifying the initial effect over the life cycle.
